%% Rechnungsvorlage – wird vom InvoiceRenderer gefüllt.
%% Platzhalter haben die Form <<NAME>> und werden vollständig ersetzt.
%% Gestaltung: TeX Gyre Heros wie im Lebenslauf, Akzentfarbe wie auf der Webseite.
\documentclass[11pt,a4paper]{article}

%% pdfx erzeugt das PDF/A-1b: XMP-Metadaten, OutputIntent mit sRGB-Profil und
%% eingebettete Schriften. Muss früh geladen werden, es zieht hyperref mit.
%% Die Metadaten kommen aus der gleichnamigen .xmpdata-Datei daneben.
\usepackage[a-1b]{pdfx}

\usepackage{fontspec}
\usepackage{geometry}
%% xcolor lädt pdfx bereits selbst – ein zweites \usepackage mit Optionen
%% gäbe einen "Option clash". colortbl direkt nachladen liefert \arrayrulecolor.
\usepackage{colortbl}
\usepackage{booktabs}  %% \midrule bringt die vertikalen Abstände selbst mit
\usepackage{tabularx}
\usepackage{array}

%% Dieselbe Schrift wie im Lebenslauf – dadurch passen die Dokumente zusammen.
\setmainfont{texgyreheros}[
  Extension      = .otf,
  UprightFont    = *-regular,
  BoldFont       = *-bold,
  ItalicFont     = *-italic,
  BoldItalicFont = *-bolditalic,
]

%% Farben von fluch-it-consulting.de
\definecolor{akzent}{HTML}{8B1A2B}
\definecolor{haupttext}{HTML}{1C1C1E}
\definecolor{gedaempft}{HTML}{5C5C62}
\definecolor{linie}{HTML}{E2E2E6}

%% Tabellenlinien global in der hellen Grautönung. Muss AUSSERHALB der Tabelle
%% stehen: Mitten im Tabellenkörper eröffnet \arrayrulecolor eine neue Zelle,
%% und das \noalign von \midrule landet dann an der falschen Stelle.
\arrayrulecolor{linie}

\geometry{a4paper, left=25mm, right=20mm, top=18mm, bottom=20mm}

\color{haupttext}
\setlength{\parindent}{0pt}
\pagestyle{empty}

%% Flattersatz statt Blocksatz. In den schmalen Spalten einer Rechnung reißt
%% Blocksatz sichtbare Löcher zwischen die Wörter; Heros verträgt Flattersatz gut.
\raggedright

%% Keine Silbentrennung: Auf einer Rechnung stehen fast nur Eigennamen,
%% Beträge und Adressen – getrennte Firmennamen sehen schlicht falsch aus.
\hyphenpenalty=10000
\exhyphenpenalty=10000
\tolerance=2000

%% Achtung: minipage ruft \@parboxrestore auf und setzt damit die
%% Absatzparameter zurück – der Flattersatz von oben gilt dort NICHT.
%% Ein \AtBeginEnvironment-Hook hilft nicht, weil \@parboxrestore erst danach
%% läuft. Deshalb steht \raggedright unten in jeder Minipage mit Fließtext.

%% Kleine, gesperrte Versalien als Rubrik – dasselbe Stilmittel wie im Lebenslauf
%% ("STANDORT", "VERFÜGBARKEIT"). \rubrik statt \label, das ist in LaTeX vergeben.
\newcommand{\rubrik}[1]{{\footnotesize\color{gedaempft}\addfontfeature{LetterSpace=14}\MakeUppercase{#1}}}

%% Dünne Trennlinie in der hellen Grautönung der Webseite.
\newcommand{\trennlinie}{\textcolor{linie}{\rule{\linewidth}{0.4pt}}}

\begin{document}

%% ------------------------------------------------------------------ Kopf
\begin{minipage}[t]{0.55\linewidth}\raggedright
  {\Large\bfseries\color{akzent}<<FIRMA>>}\\[1.5mm]
  {\small\color{gedaempft}<<INHABER>>}
\end{minipage}%
\hfill
\begin{minipage}[t]{0.40\linewidth}
  \raggedleft\small\color{gedaempft}
  <<ABSENDER_KONTAKT>>
\end{minipage}

\vspace{6mm}
\trennlinie
\vspace{10mm}

%% -------------------------------------------------- Empfänger + Metadaten
\begin{minipage}[t]{0.55\linewidth}\raggedright
  {\scriptsize\color{gedaempft}<<RUECKSENDEZEILE>>}\\[4mm]
  <<EMPFAENGER>>
\end{minipage}%
\hfill
\begin{minipage}[t]{0.40\linewidth}
  \raggedleft
  <<METADATEN>>
\end{minipage}

\vspace{14mm}

%% ------------------------------------------------------------------ Titel
{\Large\bfseries\color{akzent}Rechnung <<RECHNUNGSNUMMER>>}

\vspace{8mm}

%% --------------------------------------------------------------- Positionen
\begin{tabularx}{\linewidth}{@{}>{\raggedright\arraybackslash}X r@{\hspace{1.5mm}}l r r@{}}
  %% \mbox um jeden Spaltenkopf: In der X-Spalte ist der Inhalt sonst ein Absatz,
  %% dessen erste Zeile tiefer sitzt als die Grundlinie der r-/l-Spalten.
  \mbox{\rubrik{Beschreibung}} & \mbox{\rubrik{Menge}} & & \mbox{\rubrik{Einzelpreis}} & \mbox{\rubrik{Betrag}} \\[1.5mm]
  \midrule
<<POSITIONEN>>
\end{tabularx}

\vspace{4mm}

%% ------------------------------------------------------------------ Summen
\hfill
\begin{minipage}{0.5\linewidth}
  \begin{tabularx}{\linewidth}{@{}>{\raggedright\arraybackslash}X r@{}}
<<SUMMEN>>
  \end{tabularx}
\end{minipage}

\vspace{10mm}

%% ------------------------------------------------------- Zahlungsbedingungen
<<ZAHLUNGSHINWEIS>>

\vspace{6mm}

<<SCHLUSSSATZ>>

\vspace{4mm}

Mit freundlichen Grüßen\\[6mm]
<<INHABER>>

\vfill

%% ---------------------------------------------------------- Bankverbindung
\trennlinie
\vspace{3mm}

\begin{minipage}[t]{0.48\linewidth}\raggedright
  \rubrik{Bankverbindung}\\[2mm]
  {\small <<BANKVERBINDUNG>>}
\end{minipage}%
\hfill
\begin{minipage}[t]{0.48\linewidth}\raggedright
  \rubrik{Steuerliche Angaben}\\[2mm]
  {\small <<STEUERANGABEN>>}
\end{minipage}

\vspace{5mm}

{\centering\scriptsize\color{gedaempft}<<FUSSZEILE>>\par}

\end{document}
